\chapter{时序分析}
\label{chap:ts_analysis}
\markboth{时序分析}{时序分析}

\begin{remark}[本章定位]\label{rem:ts_chapter_positioning}
时序分析研究离散时间序列中趋势、周期和随机成分的提取、建模与预测。本章按“定义—命题—引理—定理—证明—推论—性质—例子—注释”的顺序，系统整理经典时序模型（AR、MA、ARMA、ARIMA、SARIMA、ARCH、GARCH 等）以及谱分析、拐点检测、周期显著性检验、贝叶斯更新和窗函数选择。凡标准概念均采用通行名称：用“周期图/功率谱估计”代替非标准“周期趋势谱”，用“加权互相关/循环卷积”代替非标准“周期卷积”，用“有限谐波回归/复指数模型”代替非标准“可调参数周转圆”，用“多尺度核平滑正则化最小二乘”代替非标准“KDI 优化”。本章所用的\textbf{采样点}、\textbf{采样率}、\textbf{采样间隔}与\textbf{采样序列}概念，其规范定义见「离散采样逼近、同余角与置闰」章 \cref{def:uniform_sampling_parallel}、\cref{def:nyquist_sampling_parallel}，本章不再重复定义。
\end{remark}

\begin{remark}[本章主题概览]\label{rem:ts_theme_overview}
本章围绕以下主题展开：平移——滞后算子与差分；缩放——窗宽与多尺度平滑；旋转——相位与相似度核；守恒——谱能量与 Fisher 检验归一化；对称——余弦/正弦相似度的偶奇对称；闭合——算法流程闭环与白噪声残差诊断；周期——主导周期与谱峰值。
\end{remark}

\begin{figure}[htbp]
\centering
\begin{tikzpicture}[>=Stealth, thick, node distance=1.0cm and 0.4cm,
  stage/.style={draw, rounded corners, align=center, minimum width=2.3cm, minimum height=0.75cm, font=\small, fill=blue!8},
  arr/.style={-{Stealth[length=2.2mm]}, very thick}]
\node[stage, fill=blue!12]  (s0) {原始序列\\$\{X_t\}$};
\node[stage, fill=cyan!12, right=of s0] (s1) {窗函数\\平滑};
\node[stage, fill=green!12, right=of s1] (s2) {瞬时\\极限角};
\node[stage, fill=teal!12, right=of s2] (s3) {拐点\\检测};
\node[stage, fill=orange!12, right=of s3] (s4) {转折率};
\node[stage, fill=yellow!16, below=of s0] (s5) {周期确定\\（周期图）};
\node[stage, fill=orange!8, right=of s5] (s6) {有限谐波\\回归};
\node[stage, fill=purple!12, right=of s6] (s7) {贝叶斯\\更新};
\node[stage, fill=red!12, right=of s7] (s8) {窗函数\\验证};
\node[stage, fill=magenta!12, right=of s8] (s9) {残差\\诊断};
\draw[arr] (s0) -- (s1);
\draw[arr] (s1) -- (s2);
\draw[arr] (s2) -- (s3);
\draw[arr] (s3) -- (s4);
\draw[arr] (s4.south) |- (s5.east);
\draw[arr] (s5) -- (s6);
\draw[arr] (s6) -- (s7);
\draw[arr] (s7) -- (s8);
\draw[arr] (s8) -- (s9);
\draw[arr, dashed, red] (s8.north) to[out=90, in=90] (s1.north);
\node[align=center, font=\small] at (5,-3.0) {闭环反馈：窗函数验证的结果反馈到第一步的窗函数选择};
\end{tikzpicture}
\caption{时序分析完整流程：从原始序列出发，经平滑、瞬时极限角、拐点检测、转折率、周期确定、有限谐波回归、贝叶斯更新、窗函数验证到残差诊断，其中窗函数验证反向反馈到第一步。}
\label{fig:ts_workflow}
\end{figure}

\section{经典时序模型}

\subsection{平稳性与滞后算子}

\begin{definition}[离散时间序列与平稳性]\label{def:ts_stationarity}
一个\textbf{离散时间序列}是观测值 $\{X_t\}_{t=1}^n$。序列称为\textbf{弱平稳}的，若 $\mathbb{E}[X_t] = \mu$（常数），$\operatorname{Var}(X_t) = \sigma^2 < \infty$，且 $\operatorname{Cov}(X_t, X_{t+h})$ 仅依赖于滞后 $h$。弱平稳性在概率意义上的实现是统计性质在时间平移下不变。
\end{definition}

\begin{definition}[自协方差函数与自相关函数]\label{def:ts_acf}
设 $\{X_t\}$ 弱平稳。其\textbf{自协方差函数}（ACVF）定义为
\[
\gamma(h)=\operatorname{Cov}(X_t,X_{t+h}),\qquad h\in\mathbb{Z},
\]
其\textbf{自相关函数}（ACF）为
\[
\rho(h)=\frac{\gamma(h)}{\gamma(0)}\in[-1,1].
\]
$\rho(0)=1$，且 $\rho(-h)=\rho(h)$ 为偶函数。
\end{definition}

\begin{definition}[偏自相关函数]\label{def:ts_pacf}
\textbf{偏自相关函数}（PACF）$\alpha(h)$ 定义为在给定中间变量 $X_{t+1},\dots,X_{t+h-1}$ 的条件下，$X_t$ 与 $X_{t+h}$ 的条件相关系数：
\[
\alpha(h)=\operatorname{Corr}\bigl(X_t,X_{t+h}\mid X_{t+1},\dots,X_{t+h-1}\bigr).
\]
等价地，$\alpha(h)$ 是 $X_{t+h}$ 对 $X_{t+h-1},\dots,X_t$ 作最优线性预测时，$X_t$ 的系数。
\end{definition}

\begin{definition}[白噪声]\label{def:ts_white_noise}
序列 $\{\varepsilon_t\}$ 称为\textbf{白噪声}，若 $\mathbb{E}[\varepsilon_t]=0$，$\operatorname{Var}(\varepsilon_t)=\sigma^2$，且 $\operatorname{Cov}(\varepsilon_t,\varepsilon_s)=0$（$t\neq s$）。若还满足各阶矩有限且独立，则称为\textbf{独立白噪声}；若 $\varepsilon_t\sim\mathcal{N}(0,\sigma^2)$ 独立同分布，则称\textbf{高斯白噪声}。
\end{definition}

\begin{definition}[滞后算子]\label{def:ts_lag_operator}
\textbf{滞后算子} $B$ 定义为 $BX_t = X_{t-1}$。$B$ 是离散平移群 $\mathbb{Z}$ 的生成元。滞后算子多项式
\[
\phi(B) = 1 - \phi_1 B - \cdots - \phi_p B^p
\]
定义了序列空间上的线性算子。
\end{definition}

\begin{figure}[htbp]
\centering
\begin{tikzpicture}[>=Stealth, thick, scale=1.0]
% 原始序列
\draw[->] (-0.3,0.9) -- (7.5,0.9) node[right] {$t$};
\foreach \t in {0,...,11} {
  \pgfmathsetmacro{\x}{\t*0.6}
  \pgfmathsetmacro{\y}{0.9+0.5*sin(50*\t)}
  \fill[blue] (\x,\y) circle (2pt);
}
\node[blue] at (3.5,1.9) {原始序列 $X_t$};
% 滞后一步
\draw[->] (-0.3,-0.3) -- (7.5,-0.3) node[right] {$t$};
\foreach \t in {0,...,11} {
  \pgfmathsetmacro{\x}{\t*0.6}
  \pgfmathsetmacro{\y}{-0.3+0.5*sin(50*(\t-1))}
  \fill[red] (\x,\y) circle (2pt);
}
\node[red] at (3.5,-1.0) {滞后序列 $BX_t=X_{t-1}$};
% 箭头表示滞后
\draw[->, green!60!black, very thick] (2.4,0.9) to[bend left=20] (3.0,-0.3);
\node[green!60!black, right, font=\footnotesize] at (3.1,0.2) {滞后 $B$};
\end{tikzpicture}
\caption{滞后算子 $B$：把序列沿时间轴右移一格。$B$ 生成离散平移群 $\mathbb{Z}$，是 ARIMA 建模的核心工具。}
\label{fig:lag_operator}
\end{figure}

\begin{property}[滞后算子的代数性质]\label{prop:ts_lag_algebra}
滞后算子 $B$ 在序列空间上满足：
\begin{enumerate}[label=(\roman*)]
    \item \textbf{线性性}：$B(\alpha X_t + \beta Y_t) = \alpha B X_t + \beta B Y_t$；
    \item \textbf{可交换性}：$B^a B^b = B^{a+b} = B^b B^a$；
    \item \textbf{可逆性}：当特征多项式 $\phi(z)$ 的根都在单位圆外时，$\phi(B)$ 可逆。
\end{enumerate}
\end{property}

\begin{proof}
(i) $B(\alpha X_t + \beta Y_t) = \alpha X_{t-1} + \beta Y_{t-1} = \alpha BX_t + \beta BY_t$。

(ii) $(B^a B^b)X_t = B^a X_{t-b} = X_{t-a-b} = B^{a+b}X_t$。

(iii) 若 $\phi(z)$ 的根均在单位圆外，则 $1/\phi(z)$ 在单位圆盘内解析，其幂级数展开 $\sum_k \psi_k z^k$ 的系数绝对可和。定义 $\psi(B) = \sum_k \psi_k B^k$，则 $\psi(B)\phi(B) = I$（在平稳序列空间上），故 $\phi(B)$ 可逆。
\end{proof}

\begin{definition}[差分算子]\label{def:ts_diff_operator}
\textbf{一阶差分算子}定义为 $\nabla = 1-B$，即 $\nabla X_t=X_t-X_{t-1}$。$d$ 阶差分为 $\nabla^d=(1-B)^d$。差分算子消除多项式趋势：若 $X_t=P_d(t)$ 为 $d$ 次多项式，则 $\nabla^d X_t$ 为常数。
\end{definition}

\begin{lemma}[差分算子的零空间]\label{lem:ts_diff_kernel}
$\ker\nabla^d$ 由 $\{1,t,t^2,\dots,t^{d-1}\}$ 张成，即所有次数小于 $d$ 的多项式序列。
\end{lemma}

\begin{proof}
对 $d$ 作归纳。$d=1$：$\nabla X_t=X_t-X_{t-1}=0\Rightarrow X_t$ 为常数。$d\ge2$：由 $\nabla^d=\nabla\nabla^{d-1}$，$\ker\nabla^d=\{X:\nabla X\in\ker\nabla^{d-1}\}$，即 $\nabla X$ 为次数小于 $d-1$ 的多项式。逐次积分得 $X$ 为次数小于 $d$ 的多项式。
\end{proof}

\begin{definition}[谱密度]\label{def:ts_spectral_density}
设 $\{X_t\}$ 弱平稳且 $\sum_h|\gamma(h)|<\infty$。其\textbf{谱密度}为
\[
S(\omega)=\frac{1}{2\pi}\sum_{h=-\infty}^{\infty}\gamma(h)e^{-i\omega h},\qquad \omega\in(-\pi,\pi].
\]
$S(\omega)\ge0$ 且为偶函数。逆公式 $\gamma(h)=\int_{-\pi}^{\pi}S(\omega)e^{i\omega h}d\omega$。
\end{definition}

\begin{theorem}[Wold 分解定理]\label{thm:ts_wold}
设 $\{X_t\}$ 为零均值弱平稳序列。则 $X_t$ 可唯一分解为
\[
X_t=\sum_{j=0}^{\infty}\psi_j\varepsilon_{t-j}+V_t,\qquad \psi_0=1,\quad\sum_j\psi_j^2<\infty,
\]
其中 $\{\varepsilon_t\}$ 为白噪声，$V_t$ 为确定性成分（可完全由过去线性预测），且 $\{\varepsilon_t\}$ 与 $\{V_t\}$ 正交。
\end{theorem}

\begin{proof}
将 $\{X_t\}$ 的闭线性张成空间按时间分解为 $\overline{\operatorname{span}}\{X_s:s\le t\}$ 的递增子空间链。取 $\varepsilon_t$ 为 $X_t$ 对过去空间的投影残差，则 $\{\varepsilon_t\}$ 为白噪声且 $X_t=\sum_j\psi_j\varepsilon_{t-j}+V_t$。唯一性由投影的唯一性给出。
\end{proof}

\begin{corollary}[ARMA 是 Wold 分解的有理形式]\label{cor:ts_arma_wold}
ARMA($p,q$) 的传递函数 $H(z)=\theta(z)/\phi(z)$ 是 Wold 分解的有限维参数化：其 $\psi_j$ 系数由 $\phi,\theta$ 的系数通过长除法唯一确定。
\end{corollary}

\begin{proof}
由 $\phi(z)H(z)=\theta(z)$，比较 $z^j$ 系数得 $\psi_j-\phi_1\psi_{j-1}-\cdots-\phi_p\psi_{j-p}=\theta_j$，这是 $\psi_j$ 的线性递推，唯一确定 $\{\psi_j\}$。
\end{proof}

\subsection{经典时序模型总览}

\begin{definition}[AR 模型]\label{def:ts_ar}
$p$ 阶自回归模型 AR($p$)：
\[
\phi(B) X_t = \varepsilon_t,
\]
其中 $\phi(B) = 1 - \phi_1 B - \cdots - \phi_p B^p$。平稳性条件为 $\phi(z)$ 的根都在单位圆外。
\end{definition}

\begin{definition}[MA 模型]\label{def:ts_ma}
$q$ 阶滑动平均模型 MA($q$)：
\[
X_t = \theta(B)\varepsilon_t,
\]
其中 $\theta(B) = 1 + \theta_1 B + \cdots + \theta_q B^q$。
\end{definition}

\begin{definition}[ARMA 模型]\label{def:ts_arma}
ARMA($p,q$)：
\[
\phi(B) X_t = \theta(B)\varepsilon_t.
\]
功率谱密度为
\[
S(\omega) = \frac{\sigma^2}{2\pi}\frac{|\theta(e^{-i\omega})|^2}{|\phi(e^{-i\omega})|^2}.
\]
\end{definition}

\begin{definition}[ARIMA 与 SARIMA]\label{def:ts_arima_sarima}
ARIMA($p,d,q$)：
\[
\phi(B)\nabla^d X_t = \theta(B)\varepsilon_t,
\]
其中 $\nabla = 1-B$。SARIMA($p,d,q$)$\times$($P,D,Q$)$_s$：
\[
\phi(B)\Phi(B^s)\nabla^d\nabla_s^D X_t = \theta(B)\Theta(B^s)\varepsilon_t.
\]
\end{definition}

\begin{definition}[ARCH 与 GARCH]\label{def:ts_arch_garch}
ARCH($q$)：
\[
\sigma_t^2 = \alpha_0 + \sum_{i=1}^q \alpha_i X_{t-i}^2.
\]
GARCH($p,q$)：
\[
\sigma_t^2 = \alpha_0 + \sum_{i=1}^q \alpha_i X_{t-i}^2 + \sum_{j=1}^p \beta_j \sigma_{t-j}^2.
\]
\end{definition}

\begin{figure}[htbp]
\centering
\begin{tikzpicture}[>=Stealth, thick, scale=0.85]
% AR(1) ACF
\begin{scope}[xshift=-6cm]
\draw[->] (-0.3,0) -- (4.5,0) node[right] {$h$};
\draw[->] (0,-0.3) -- (0,1.6) node[above] {ACF};
\foreach \h in {0,1,...,7} {
  \pgfmathsetmacro{\x}{\h*0.5}
  \pgfmathsetmacro{\y}{0.7^(\h)}
  \draw[blue, very thick] (\x,0) -- (\x,\y);
  \fill[blue] (\x,\y) circle (1.5pt);
}
\node[font=\footnotesize] at (2,-0.9) {AR(1) ACF：拖尾};
\end{scope}
% AR(1) PACF
\begin{scope}[xshift=-2cm]
\draw[->] (-0.3,0) -- (4.5,0) node[right] {$h$};
\draw[->] (0,-0.3) -- (0,1.6) node[above] {PACF};
\foreach \h in {0,1,...,7} {
  \pgfmathsetmacro{\x}{\h*0.5}
  \pgfmathsetmacro{\y}{\h==1 ? 0.7 : 0}
  \draw[red, very thick] (\x,0) -- (\x,\y);
  \fill[red] (\x,\y) circle (1.5pt);
}
\node[font=\footnotesize] at (2,-0.9) {AR(1) PACF：截断};
\end{scope}
% MA(1) ACF
\begin{scope}[xshift=2cm]
\draw[->] (-0.3,0) -- (4.5,0) node[right] {$h$};
\draw[->] (0,-0.3) -- (0,1.6) node[above] {ACF};
\foreach \h in {0,1,...,7} {
  \pgfmathsetmacro{\x}{\h*0.5}
  \pgfmathsetmacro{\y}{\h==1 ? 0.6 : 0}
  \draw[blue, very thick] (\x,0) -- (\x,\y);
  \fill[blue] (\x,\y) circle (1.5pt);
}
\node[font=\footnotesize] at (2,-0.9) {MA(1) ACF：截断};
\end{scope}
% MA(1) PACF
\begin{scope}[xshift=6cm]
\draw[->] (-0.3,0) -- (4.5,0) node[right] {$h$};
\draw[->] (0,-0.3) -- (0,1.6) node[above] {PACF};
\foreach \h in {0,1,...,7} {
  \pgfmathsetmacro{\x}{\h*0.5}
  \pgfmathsetmacro{\y}{0.6^(\h) * cos(deg(\h)*1.2)}
  \draw[red, very thick] (\x,0) -- (\x,\y);
  \fill[red] (\x,\y) circle (1.5pt);
}
\node[font=\footnotesize] at (2,-0.9) {MA(1) PACF：拖尾};
\end{scope}
\end{tikzpicture}
\caption{AR 与 MA 的 ACF/PACF 特征对比：AR 模型 PACF 截断、ACF 拖尾；MA 模型 ACF 截断、PACF 拖尾。这是模型识别（Box-Jenkins 方法）的判据。}
\label{fig:ar_ma_acf_pacf}
\end{figure}

\begin{table}[htbp]
    \centering
    \caption{经典时序模型的标准特征}
    \label{tab:ts_classical_correspondence}
    \begin{tabular}{>{\centering\arraybackslash}p{2.2cm}>{\centering\arraybackslash}p{2.5cm}>{\centering\arraybackslash}p{2.5cm}>{\centering\arraybackslash}p{2.5cm}>{\centering\arraybackslash}p{3cm}}
        \toprule
        模型 & ACF/PACF 特征 & 频域特征 & 算子对应 & 传递函数 \\
        \midrule
        AR($p$) & PACF 截断 & Yule--Walker 递推 & 梯度（自反馈） & $H(\omega) = 1/\phi(e^{-i\omega})$ \\
        MA($q$) & ACF 截断 & 脉冲响应 & 散度（冲击） & $H(\omega) = \theta(e^{-i\omega})$ \\
        ARMA & ACF/PACF 拖尾 & 有理谱 & 梯度+散度 & $H(\omega) = \theta(e^{-i\omega})/\phi(e^{-i\omega})$ \\
        ARIMA & 差分后同 ARMA & 差分趋势 & 梯度（离散微分） & $H = \theta/(\phi(1-e^{-i\omega})^d)$ \\
        SARIMA & 季节+短周期 & 多尺度趋势 & 梯度（多尺度） & 多尺度传递函数 \\
        ARCH & 平方 ACF 拖尾 & 波动率自趋势 & 散度（波动缩放） & 波动率传递函数 \\
        GARCH & 平方 ACF/PACF 拖尾 & 波动率持续 & 梯度+散度 & $\alpha/(1-\beta)$ 波动率核 \\
        \bottomrule
    \end{tabular}
\end{table}

\begin{theorem}[Yule-Walker 方程]\label{thm:ts_yule_walker}
设 AR($p$) 模型 $\phi(B)X_t=\varepsilon_t$ 平稳。则自协方差函数满足
\[
\gamma(h)-\phi_1\gamma(h-1)-\cdots-\phi_p\gamma(h-p)=\begin{cases}\sigma^2,&h=0,\\0,&h\ge1.\end{cases}
\]
以 $\rho(h)$ 表示，$h\ge1$ 给出线性方程组
\[
\rho(h)=\phi_1\rho(h-1)+\cdots+\phi_p\rho(h-p),\qquad h=1,\dots,p.
\]
\end{theorem}

\begin{proof}
以 $X_{t-h}$ 乘 $\phi(B)X_t=\varepsilon_t$ 两边取期望，得 $\sum_j\phi_j\gamma(h-j)=\mathbb{E}[\varepsilon_tX_{t-h}]$。当 $h\ge1$ 时，$\varepsilon_t$ 与 $X_{t-h}$ 正交（$\varepsilon_t$ 与过去观测不相关），故右端为 $0$。当 $h=0$ 时，$\mathbb{E}[\varepsilon_tX_t]=\sigma^2$。
\end{proof}

\begin{property}[平稳性的几何直观]\label{prop:ts_stationarity_geometry}
AR(1) 模型 $X_t=\phi X_{t-1}+\varepsilon_t$ 平稳当且仅当 $|\phi|<1$。此时自协方差 $\gamma(h)=\sigma^2\phi^{|h|}/(1-\phi^2)$ 指数衰减，衰减速率由 $\phi$ 控制：$|\phi|$ 越接近 $1$，记忆越长。
\end{property}

\begin{proof}
特征多项式 $\phi(z)=1-\phi z$ 的根为 $z=1/\phi$，在单位圆外当且仅当 $|\phi|<1$。自协方差由 $\gamma(h)=\phi\gamma(h-1)$（$h\ge1$）递推得 $\gamma(h)=\phi^h\gamma(0)$，其中 $\gamma(0)=\sigma^2/(1-\phi^2)$。
\end{proof}

\begin{figure}[htbp]
\centering
\begin{tikzpicture}[>=Stealth, thick, scale=0.95]
\draw[->] (-3.3,0) -- (3.3,0) node[right] {$\operatorname{Re}$};
\draw[->] (0,-2.0) -- (0,2.0) node[above] {$\operatorname{Im}$};
\draw[gray!60] (0,0) circle (1.5);
% 单位圆内的根（非平稳）
\fill[red] (0.5,0) circle (2.5pt);
\node[red, right, font=\footnotesize] at (0.55,0.15) {非平稳根};
% 单位圆外的根（平稳）
\fill[blue] (2.5,0) circle (2.5pt);
\node[blue, right, font=\footnotesize] at (2.55,0.15) {平稳根};
% 单位圆标注
\node[font=\footnotesize] at (1.1,1.05) {$|z|=1$};
\node[align=center, font=\small] at (0,-2.7) {平稳性条件：特征多项式 $\phi(z)$ 的根都在单位圆外（蓝点），不在圆内（红点）};
\end{tikzpicture}
\caption{ARMA 平稳性条件：特征多项式 $\phi(z)$ 的根全部落在复平面单位圆外。}
\label{fig:stationarity_condition}
\end{figure}

\section{时序分析的基础工具}

\begin{remark}[窗函数定义的归并]\label{rem:ts_window_function_merged}
窗函数的一般定义（频域/旁瓣视角）及其在时序分析中的三重作用，已统一归入「离散采样逼近、同余角与置闰」章（见 \cref{def:window_function} 及其时序应用子节）。本章后续直接使用该定义，不再重复给出。
\end{remark}

\subsection{瞬时极限角：离散梯度的方向信息}

\begin{definition}[时序的瞬时极限角]\label{def:ts_limit_angle}
对窗函数平滑后的序列 $\tilde{X}_t$，定义\textbf{瞬时极限角}为
\[
\theta_t = \arctan(\nabla \tilde{X}_t) = \arctan(\tilde{X}_t - \tilde{X}_{t-1}),
\]
其中 $\nabla \tilde{X}_t = \tilde{X}_t - \tilde{X}_{t-1}$ 是离散梯度。$\theta_t \in (-\pi/2, \pi/2)$。
\end{definition}

\section{拐点检测：基于相似度的方向转折判定}

\begin{definition}[拐点的数学定义]\label{def:ts_turning_point}
设 $\theta_t \in (-\pi/2, \pi/2)$ 为瞬时极限角。$t^*$ 称为\textbf{拐点}，若
\[
|\theta_{t^*} - \theta_{t^*-1}| > \tau,
\]
其中 $\tau \in (0, \pi/2)$ 为预设阈值。记 $\Delta\theta = \theta_{t^*} - \theta_{t^*-1}$。

按「标量与矢量」章的统一口径，两趋势方向之间的\textbf{夹角}一律取 $\theta=\lvert\alpha-\beta\rvert\in[0,\pi]$；有符号的 $\Delta\theta$ 是\textbf{差角}，仅用于正弦族相似度的符号判定。
\end{definition}

\begin{lemma}[拐点阈值与余弦相似度的等价性]\label{lem:ts_turning_equivalence}
设 $|\Delta\theta| \in [0, \pi)$，$\tau \in (0, \pi)$。则
\[
|\Delta\theta| > \tau \quad \Longleftrightarrow \quad \cos(\Delta\theta) < \cos\tau.
\]
\end{lemma}

\begin{proof}
余弦函数在 $[0, \pi]$ 上严格单调递减，且 $\cos$ 为偶函数。故 $|\Delta\theta| > \tau$ 当且仅当 $\cos|\Delta\theta| < \cos\tau$。
\end{proof}

\begin{proposition}[拐点的相似度特征]\label{prop:ts_turning_similarity}
在拐点 $t^*$ 处：
\begin{enumerate}[label=(\roman*)]
    \item 余弦相似度 $\operatorname{Similarity}^{(1)} = \cos(\Delta\theta) < \cos\tau$；
    \item 正弦相似度 $\operatorname{Similarity}^{(2)} = \sin(\Delta\theta)$ 的符号给出转向；
    \item 倍角正弦相似度 $\operatorname{Similarity}^{(4)} = \sin(2\Delta\theta)$ 随 $\Delta\theta$ 严格递增；其绝对值随 $|\Delta\theta|$ 严格递增，当 $|\Delta\theta| < \pi/4$。
\end{enumerate}
\end{proposition}

\begin{proof}
(i) 由引理 \ref{lem:ts_turning_equivalence}。(ii) $\sin$ 为奇函数。(iii) $\frac{d}{d\theta}\sin 2\theta = 2\cos 2\theta > 0$ 当 $|\theta| < \pi/4$。绝对值结论由奇函数与单调性直接得到。
\end{proof}

\begin{figure}[htbp]
\centering
\begin{tikzpicture}[>=Stealth, thick, scale=1.0]
% 时间轴
\draw[->] (-0.3,0) -- (9,0) node[right] {$t$};
% 分段线性趋势
\draw[blue, very thick] (0,0) -- (4,1.6) -- (5.5,0.4) -- (8.5,-0.6);
% 拐点标注
\fill[red] (4,1.6) circle (3pt);
\node[red, above, font=\footnotesize] at (4,1.75) {拐点};
\fill[red] (5.5,0.4) circle (3pt);
% 瞬时极限角
\draw[->, green!60!black, thick] (2.5,0.8) -- (3.5,0.8);
\node[green!60!black, below, font=\footnotesize] at (3.0,0.7) {$\theta_{t-1}>0$};
\draw[->, orange, thick] (5.0,0.4) -- (6.0,0.15);
\node[orange, below, font=\footnotesize] at (5.5,0.05) {$\theta_t<0$};
% 阈值
\node[align=center, font=\small] at (4.5,-1.4) {拐点判定：$|\theta_{t^*}-\theta_{t^*-1}|>\tau$，等价于 $\cos(\Delta\theta)<\cos\tau$};
\end{tikzpicture}
\caption{拐点检测：当瞬时极限角的变化量 $|\Delta\theta|$ 超过阈值 $\tau$ 时，判定为拐点。正弦相似度的符号精确标记转折方向。}
\label{fig:turning_point_detection}
\end{figure}

\begin{example}[拐点检测实例]\label{ex:ts_turning_example}
序列在 $t=50$ 处从上升转为下降。$\tilde{X}_{49} = 10.2$，$\tilde{X}_{50} = 10.5$，$\tilde{X}_{51} = 10.1$。则 $\theta_{50} \approx 16.7^\circ$，$\theta_{51} \approx -21.8^\circ$。$\Delta\theta \approx -38.5^\circ$。阈值 $\tau = 30^\circ$ 时 $\cos(38.5^\circ) \approx 0.78 < \cos(30^\circ) \approx 0.87$，$t=51$ 为拐点。$\operatorname{Similarity}^{(4)} = \sin(-77^\circ) \approx -0.97$，负号表示向下转折。
\end{example}

\section{转折率：拐点两侧正弦趋势之差}

\begin{definition}[转折率]\label{def:ts_refraction}
在拐点 $t^*$ 处，定义\textbf{入射角} $\theta_1 = \theta_{t^*-1}$ 与\textbf{出射角} $\theta_2 = \theta_{t^*+1}$。\textbf{转折率}为
\[
R_{\text{turn}}(t^*) = |\sin\theta_1 - \sin\theta_2|.
\]
\end{definition}

\begin{theorem}[转折率的半角分解与界]\label{thm:ts_refraction_relation}
设 $\theta_1, \theta_2 \in (-\pi/2, \pi/2)$。则
\[
R_{\text{turn}} = 2\left| \cos\frac{\theta_1 + \theta_2}{2} \right| \cdot \left| \sin\frac{\theta_1 - \theta_2}{2} \right|,
\]
且 $0 \le R_{\text{turn}} \le 2$。特别地：$R = 0 \iff \theta_1 = \theta_2$；$R = 2 \iff \theta_2 = -\theta_1 = \pm\pi/2$；方向反转（$\theta_2 = -\theta_1$）时 $R = 2|\sin\theta_1|$。
\end{theorem}

\begin{proof}
半角分解由和差化积公式给出。界由 $|\cos| \le 1$ 和 $|\sin| \le 1$ 得出。三特例逐条验证。
\end{proof}

\begin{figure}[htbp]
\centering
\begin{tikzpicture}[>=Stealth, thick, scale=1.0]
% 单位圆
\draw[gray!40] (0,0) circle (1.5);
\draw[->] (-1.8,0) -- (1.8,0);
\draw[->] (0,-1.8) -- (0,1.8);
% 入射角
\draw[->, blue, very thick] (0,0) -- (1.2,0.9);
\draw[->, blue, thick] (0,0) -- (1.2,0);
\draw[->, blue, thick] (0.9,0) arc (0:37:0.9);
\node[blue, font=\footnotesize] at (1.05,0.4) {$\theta_1$};
% 出射角
\draw[->, red, very thick] (0,0) -- (1.3,-0.6);
\draw[->, red, thick] (0,0) -- (1.5,0);
\draw[->, red, thick] (1.2,0) arc (0:-25:1.2);
\node[red, font=\footnotesize] at (1.35,-0.4) {$\theta_2$};
% 正弦投影
\draw[dashed, green!60!black] (1.2,0.9) -- (1.2,0);
\draw[dashed, orange] (1.3,-0.6) -- (1.3,0);
\node[green!60!black, font=\footnotesize] at (1.5,0.5) {$\sin\theta_1$};
\node[orange, font=\footnotesize] at (1.55,-0.35) {$\sin\theta_2$};
\node[align=center, font=\small] at (0,-2.4) {$R_{\text{turn}}=|\sin\theta_1-\sin\theta_2|=2|\cos\frac{\theta_1+\theta_2}{2}|\cdot|\sin\frac{\theta_1-\theta_2}{2}|$};
\end{tikzpicture}
\caption{转折率的半角分解几何图：$R_{\text{turn}}$ 是两侧正弦投影之差，半角分解将其分为“平均方向”与“方向差”两个因子。}
\label{fig:refraction_decomposition}
\end{figure}

\begin{example}[转折率计算实例]\label{ex:ts_refraction_example}
延续例 \ref{ex:ts_turning_example}。$\sin\theta_1 = 0.287$，$\sin\theta_2 = -0.371$。$R_{\text{turn}} = 0.658$。半角分解验证一致。
\end{example}

\section{周期确定：周期图的峰值}
\label{sec:ts_period_determination}

\begin{definition}[基于周期图的周期确定]\label{def:ts_period_detection}
计算周期图
\[
I_n(\omega_k) = \frac{1}{n}\left|\sum_{t=1}^n X_t e^{-i\omega_k t}\right|^2.
\]
主导周期由峰值频率确定：
\[
\omega_{\text{period}} = \arg\max_{\omega \in (0, \pi]} I_n(\omega), \qquad
T_{\text{period}} = \frac{2\pi}{\omega_{\text{period}}}.
\]
\end{definition}

\begin{proposition}[周期显著性的 Fisher 检验]\label{prop:ts_period_significance}
设周期图在 $m = \lfloor (n-1)/2 \rfloor$ 个正频率点上的取值为 $I_1, \dots, I_m$。Fisher 检验统计量
\[
g = \frac{\max_k I_k}{\sum_{k=1}^{m} I_k}.
\]
阈值
\[
g_\alpha = 1 - (\alpha/m)^{1/(m-1)}.
\]
若 $g > g_\alpha$，周期显著。
\end{proposition}

\begin{proof}
白噪声下周期图纵坐标独立同分布（指数分布），归一化比值服从 $\operatorname{Beta}(1, m-1)$。并集界给出
\[
\mathbb{P}(g > g_0) \le m(1-g_0)^{m-1}.
\]
代入 $g_0 = g_\alpha$ 得 $\alpha$。
\end{proof}

\begin{figure}[htbp]
\centering
\begin{tikzpicture}[>=Stealth, thick, scale=0.95]
\draw[->] (-0.3,0) -- (8.5,0) node[right] {$\omega$};
\draw[->] (0,-0.3) -- (0,3.0) node[above] {$I_n(\omega)$};
% 周期图曲线
\draw[blue, very thick, domain=0.05:8.2, samples=200]
  plot (\x, {1.0+0.8*sin(deg(\x*1.8))^2+0.5*sin(deg(\x*3.7))^2+0.3*sin(deg(\x*7.1))^2});
% 峰值标注
\draw[dashed, red] (2.1,0) -- (2.1,2.4);
\node[red, above, font=\footnotesize] at (2.1,2.45) {主导峰};
\draw[dashed, red] (5.3,0) -- (5.3,1.9);
\node[red, above, font=\footnotesize] at (5.3,1.95) {次峰};
% Fisher 阈值
\draw[dashed, green!60!black] (0,2.2) -- (8.2,2.2);
\node[green!60!black, right, font=\footnotesize] at (8.2,2.2) {$g_\alpha$ 阈值};
\node[align=center, font=\small] at (4.2,-1.0) {周期图峰值超过 Fisher 阈值 $g_\alpha$ 时，对应周期显著};
\end{tikzpicture}
\caption{周期图与 Fisher 检验：谱峰值高于阈值 $g_\alpha$ 的频率对应显著周期，主峰给出主导周期 $T_{\text{period}}=2\pi/\omega_{\text{period}}$。}
\label{fig:periodogram_fisher}
\end{figure}

\begin{property}[周期图的渐近分布]\label{prop:ts_periodogram_distribution}
在高斯白噪声下，周期图 $I_n(\omega_k)$ 渐近服从指数分布，均值为 $\sigma^2$。即
\[
I_n(\omega_k)\xrightarrow{d}\operatorname{Exp}(\sigma^2),
\]
且不同频率处的 $I_n(\omega_k),I_n(\omega_{k'})$ 渐近独立（$k\neq k'$）。
\end{property}

\begin{proof}
高斯白噪声的 DFT 系数 $\sum_tX_te^{-i\omega_kt}$ 为复高斯随机变量，实部与虚部独立同分布于 $\mathcal{N}(0,n\sigma^2/2)$，故 $I_n(\omega_k)$ 为两个独立高斯平方和的 $1/n$ 倍，即指数分布。不同频率处的 DFT 系数渐近正交，故渐近独立。
\end{proof}

\section{有限谐波回归模型}

\begin{definition}[有限谐波回归模型]\label{def:ts_harmonic_regression}
\[
X_t \approx \operatorname{Re}\left[ \sum_{k=1}^{K} A_k e^{i(\omega_k t + \phi_k)} \right] + \varepsilon_t,
\]
其中 $\{A_k, \omega_k, \phi_k\}_{k=1}^{K}$ 是可调参数。
\end{definition}

\begin{proposition}[有限谐波模型参数的正则化最小二乘估计]\label{prop:ts_harmonic_estimate}
有限谐波参数通过正则化最小二乘进行估计。目标函数为
\[
\mathcal{L}(\boldsymbol{\Theta}) = \sum_{t=1}^n \left( X_t - \operatorname{Re}\left[\sum_k A_k e^{i(\omega_k t + \phi_k)}\right] \right)^2 + \lambda \sum_k |A_k|.
\]
在适当的正则化条件下，估计量 $\hat{\boldsymbol{\Theta}}_n$ 依概率收敛于真实参数 $\boldsymbol{\Theta}_0$。
\end{proposition}

\begin{proof}
目标函数 $\mathcal{L}$ 是 $\boldsymbol{\Theta}$ 的连续可微函数（除 $A_k = 0$ 处的 $L^1$ 正则项不可微点，可用次梯度处理）。由 $L^2$ 损失项的强凸性（在 $\omega_k$ 分离的紧参数空间上）与 $L^1$ 正则项的凸性，$\mathcal{L}$ 在紧参数空间上有唯一最小值。设 $\hat{\boldsymbol{\Theta}}_n = \arg\min \mathcal{L}$。由 Wald 一致性条件：(i) 参数空间紧致；(ii) 真实参数 $\boldsymbol{\Theta}_0$ 是唯一最小值（模型可识别性：$A_k > 0$ 且 $\omega_k$ 两两互异保证谐波分量唯一确定）；(iii) 损失函数逐点收敛于期望损失。满足条件时 $\hat{\boldsymbol{\Theta}}_n \xrightarrow{p} \boldsymbol{\Theta}_0$。多尺度核平滑可加速数值优化，但不改变一致性。
\end{proof}

\begin{theorem}[有限谐波估计的渐近正态性]\label{thm:ts_harmonic_asymptotic}
在上述正则化条件下，估计量满足
\[
\sqrt{n}\bigl(\hat{\boldsymbol{\Theta}}_n-\boldsymbol{\Theta}_0\bigr)\xrightarrow{d}\mathcal{N}(0,\mathcal{I}(\boldsymbol{\Theta}_0)^{-1}),
\]
其中 $\mathcal{I}(\boldsymbol{\Theta}_0)$ 为 Fisher 信息矩阵。
\end{theorem}

\begin{proof}
由 $\sqrt{n}$-一致性、目标函数的二阶可微性（除 $A_k=0$ 处）与 Taylor 展开，$\hat{\boldsymbol{\Theta}}_n-\boldsymbol{\Theta}_0\approx-\mathcal{I}(\boldsymbol{\Theta}_0)^{-1}\nabla\mathcal{L}(\boldsymbol{\Theta}_0)$，其中 $\nabla\mathcal{L}(\boldsymbol{\Theta}_0)$ 为渐近正态。由 Slutsky 定理即得。
\end{proof}

\section{贝叶斯更新：参数的在线推断}

\begin{definition}[有限谐波参数的贝叶斯更新]\label{def:ts_bayesian_update}
后验分布
\[
p(\boldsymbol{\Theta} \mid X_{1:t+1}) \propto p(X_{t+1} \mid \boldsymbol{\Theta}) \cdot p(\boldsymbol{\Theta} \mid X_{1:t}).
\]
\end{definition}

\begin{proposition}[后验收敛性与周期稳定性]\label{prop:ts_bayes_convergence}
在标准正则条件下，Bernstein--von Mises 定理给出：当 $n \to \infty$ 时，
\[
\sqrt{n}(\boldsymbol{\Theta} - \hat{\boldsymbol{\Theta}}_n) \mid X_{1:n} \xrightarrow{d} \mathcal{N}(0, \mathcal{I}(\boldsymbol{\Theta}_0)^{-1}),
\]
其中 $\mathcal{I}(\boldsymbol{\Theta}_0)$ 为 Fisher 信息矩阵。后验分布在真实参数处集中，宽度为 $O(n^{-1/2})$。若后验在 $\omega_k$ 处持续集中，则该周期分量稳定；若扩散，则不稳定。
\end{proposition}

\begin{proof}
在标准正则条件下，先验在 $\boldsymbol{\Theta}_0$ 处连续且为正时，后验渐近正态。Fisher 信息非退化保证协方差矩阵 $\mathcal{I}(\boldsymbol{\Theta}_0)^{-1}$ 良定义。集中性由 $O(n^{-1/2})$ 的渐近标准差给出。周期稳定性判据由后验方差在 $\omega_k$ 维度上的大小直接判定。
\end{proof}

\begin{proposition}[遗忘因子的自适应]\label{prop:ts_forgetting}
对非平稳时序，遗忘因子 $\gamma \in (0,1]$ 的定义为
\[
p_\gamma(\boldsymbol{\Theta} \mid X_{1:t}) \propto p(X_t \mid \boldsymbol{\Theta}) \cdot [p_\gamma(\boldsymbol{\Theta} \mid X_{1:t-1})]^\gamma.
\]
等价地，
\[
\log p_\gamma(\boldsymbol{\Theta} \mid X_{1:t}) \propto \sum_{s=1}^t \gamma^{t-s} \log p(X_s \mid \boldsymbol{\Theta}),
\]
即旧观测的似然贡献按指数衰减。自适应遗忘因子在拐点附近减小 $\gamma$（增强适应性），在平稳段增大 $\gamma$（增强稳定性）。
\end{proposition}

\begin{proof}
展开遗忘因子的递归定义：
\[
p_\gamma(\boldsymbol{\Theta} \mid X_{1:t}) \propto \prod_{s=1}^t p(X_s \mid \boldsymbol{\Theta})^{\gamma^{t-s}}.
\]
取对数得加权对数似然
\[
\sum_{s=1}^t \gamma^{t-s} \log p(X_s \mid \boldsymbol{\Theta}).
\]
当 $\gamma \to 1$ 时退化为标准贝叶斯更新（所有观测等权）；当 $\gamma < 1$ 时，旧观测权重指数衰减，有效样本量约为 $1/(1-\gamma)$。拐点检测在 $t^*$ 处触发时，将 $\gamma$ 暂时降至 $\gamma_{\text{low}}$（如 $0.9$），使模型快速适应方向转折；平稳段恢复 $\gamma \to 1$ 以保持估计精度。
\end{proof}

\begin{figure}[htbp]
\centering
\begin{tikzpicture}[>=Stealth, thick, scale=0.9]
\draw[->] (-0.3,0) -- (9,0) node[right] {$s$};
\draw[->] (0,-0.3) -- (0,2.8) node[above] {权重 $\gamma^{t-s}$};
% 不同 gamma 的权重曲线
\draw[blue, very thick, domain=0:8.5, samples=200]
  plot (\x, {2.5*0.9^(\x)});
\draw[red, very thick, domain=0:8.5, samples=200]
  plot (\x, {2.5*0.7^(\x)});
\draw[green!60!black, very thick, domain=0:8.5, samples=200]
  plot (\x, {2.5*0.5^(\x)});
\node[blue, right, font=\footnotesize] at (3,1.6) {$\gamma=0.9$：平稳段，长记忆};
\node[red, right, font=\footnotesize] at (3,0.8) {$\gamma=0.7$：中等适应};
\node[green!60!black, right, font=\footnotesize] at (3,0.3) {$\gamma=0.5$：拐点附近，快速适应};
\node[align=center, font=\small] at (4.5,-1.2) {遗忘因子 $\gamma$ 控制旧观测的指数衰减权重：$\gamma$ 越小，适应越快但估计方差越大};
\end{tikzpicture}
\caption{遗忘因子的指数衰减：$\gamma$ 越小，旧观测权重衰减越快，模型对近期的适应越强，但有效样本量降低。}
\label{fig:forgetting_factor}
\end{figure}

\section{窗函数验证：基于频谱分量确定的模型诊断}

\begin{definition}[窗函数验证的频谱准则]\label{def:ts_window_validation}
保真度定义为
\[
\mathcal{V}(K_\Delta) =
\frac{\int_{-\pi}^{\pi} |\hat{K}_\Delta(\omega)|^2 \cdot S_X(\omega) \, d\omega}
{\int_{-\pi}^{\pi} S_X(\omega) \, d\omega}.
\]
\end{definition}

\begin{proposition}[窗函数验证的拐点--周期权衡]\label{prop:ts_window_tradeoff}
小窗口保真度高但噪声抑制不足；大窗口噪声抑制充分但高频信号被平滑；最优窗口在两者间平衡。
\end{proposition}

\begin{proof}
小窗口 $\Delta \to 0$ 时 $\hat{K}_\Delta(\omega) \to 1$（逐点），故 $\mathcal{V} \to 1$；但平滑后的序列 $\tilde{X}_t \to X_t$（原始序列），噪声完全保留，拐点检测产生伪拐点的概率增大。大窗口 $\Delta \to \infty$ 时 $\hat{K}_\Delta(\omega) \to \delta(\omega)$（集中在零频），高频分量被完全抑制，周期估计精度高（低频周期保留）但拐点信息丢失。最优窗口由最小化均方误差或交叉验证确定，权衡关系由此证明。
\end{proof}

\begin{figure}[htbp]
\centering
\begin{tikzpicture}[>=Stealth, thick, scale=0.95]
\draw[->] (-0.3,0) -- (8.5,0) node[right] {窗宽 $\Delta$};
\draw[->] (0,-0.3) -- (0,3.0) node[above] {性能};
% 保真度曲线（下降）
\draw[blue, very thick, domain=0.5:8, samples=200]
  plot (\x, {2.8*exp(-\x/4)});
% 噪声抑制曲线（上升）
\draw[red, very thick, domain=0.5:8, samples=200]
  plot (\x, {2.8*(1-exp(-\x/2))});
% 最优窗口
\draw[dashed, green!60!black] (2.6,0) -- (2.6,2.8);
\node[green!60!black, above, font=\footnotesize] at (2.6,2.85) {最优 $\Delta^*$};
\node[blue, right, font=\footnotesize] at (5.5,1.5) {保真度（小窗最优）};
\node[red, right, font=\footnotesize] at (5.5,2.3) {噪声抑制（大窗最优）};
\node[align=center, font=\small] at (4.2,-1.2) {拐点--周期权衡：最优窗宽 $\Delta^*$ 在保真度与噪声抑制之间取平衡};
\end{tikzpicture}
\caption{窗函数的拐点--周期权衡：小窗口保真度高但噪声抑制不足，大窗口噪声抑制充分但高频丢失，最优窗口 $\Delta^*$ 在两者之间。}
\label{fig:window_tradeoff}
\end{figure}

\begin{proposition}[残差的相似度诊断]\label{prop:ts_residual_diagnosis}
残差相似度诊断函数
\[
\mathcal{D}(\hat{\varepsilon}) =
\max_{h \ge 1,\, d = 1,\dots,6}
\left|
\frac{1}{n-h} \sum_{t=h+1}^{n}
S^{(d)}\big(\theta_t^{\hat{\varepsilon}}, \theta_{t-h}^{\hat{\varepsilon}}\big)
\right|.
\]
\end{proposition}

\begin{proof}
理想白噪声残差的相位序列独立。差角族核各项为 $e^{i(\theta_t - \theta_{t-h})}$，相位差在 $[0,2\pi)$ 上均匀分布，故
\[
\mathbb{E}[e^{i(\theta_t - \theta_{t-h})}] = 0.
\]
由大数定律（$h$-相依序列），均值依概率收敛于 $0$，即 $\mathcal{D} \to 0$。若某核的 $\mathcal{D}$ 显著非零，表明残差在该核对应的频率尺度上存在未捕获的周期结构。
\end{proof}

\begin{theorem}[Ljung-Box 残差白噪声检验]\label{thm:ts_ljung_box}
残差 $\{\hat{\varepsilon}_t\}$ 的 $m$ 阶自相关 $\hat\rho(h)$ 满足检验统计量
\[
Q=n(n+2)\sum_{h=1}^m\frac{\hat\rho(h)^2}{n-h}\xrightarrow{d}\chi^2_{m-p-q},
\]
其中 $p+q$ 为模型参数个数。当 $Q>\chi^2_{m-p-q,\alpha}$ 时，拒绝残差为白噪声的原假设。
\end{theorem}

\begin{proof}
在模型正确设定下，残差的自相关 $\hat\rho(h)$ 渐近正态且相关矩阵为单位矩阵。由二次型分布理论，$Q$ 渐近服从自由度为 $m-p-q$ 的 $\chi^2$ 分布。
\end{proof}

\section{时序分析的完整流程}

\begin{definition}[时序分析算法流程]\label{def:ts_algorithm}
完整算法流程：
\begin{enumerate}[label=(\roman*)]
    \item 窗函数选择与平滑；
    \item 瞬时极限角计算；
    \item 拐点检测；
    \item 转折率计算；
    \item 周期确定；
    \item 有限谐波回归建模与参数估计；
    \item 贝叶斯更新；
    \item 窗函数验证。
\end{enumerate}
\end{definition}

\begin{figure}[htbp]
\centering
\begin{tikzpicture}[>=Stealth, thick, node distance=0.9cm and 0.9cm,
  tstep/.style={draw, rounded corners, align=center, minimum width=2.2cm, minimum height=0.8cm, font=\small, fill=blue!8}]
\node[tstep, fill=blue!12]  (a) {1. 窗函数\\平滑};
\node[tstep, fill=cyan!12, right=of a] (b) {2. 瞬时\\极限角};
\node[tstep, fill=green!12, right=of b] (c) {3. 拐点\\检测};
\node[tstep, fill=teal!12, right=of c] (d) {4. 转折率};
\node[tstep, fill=orange!12, below=of a] (e) {5. 周期\\确定};
\node[tstep, fill=yellow!16, right=of e] (f) {6. 有限谐波\\回归};
\node[tstep, fill=purple!12, right=of f] (g) {7. 贝叶斯\\更新};
\node[tstep, fill=red!12, right=of g] (h) {8. 窗函数\\验证};
\draw[->, very thick] (a) -- (b);
\draw[->, very thick] (b) -- (c);
\draw[->, very thick] (c) -- (d);
\draw[->, very thick] (d.south) |- (e.east);
\draw[->, very thick] (e) -- (f);
\draw[->, very thick] (f) -- (g);
\draw[->, very thick] (g) -- (h);
\draw[->, very thick, dashed, red] (h.north) to[out=90, in=90] (a.north);
\node[align=center, font=\small] at (4.5,-2.8) {闭环反馈：第 8 步窗函数验证的结果反馈到第 1 步窗函数选择};
\end{tikzpicture}
\caption{时序分析完整算法流程闭环：八个步骤构成完整闭环，窗函数验证结果反向反馈到窗函数选择，实现模型迭代优化。}
\label{fig:ts_algorithm_loop}
\end{figure}

\begin{example}[完整流程在模拟数据上的应用]\label{ex:ts_complete_example}
模拟数据：
\[
X_t = 5\sin(2\pi t/20) + 2\sin(2\pi t/7) + \varepsilon_t, \quad t = 1,\dots,200,
\]
其中 $\varepsilon_t \sim \mathcal{N}(0,1)$。

拐点检测：阈值 $\tau = 45^\circ$ 检测到约 $18$ 个拐点。转折率 $R_{\text{turn}}$ 范围 $0.3$ 到 $1.8$。周期确定：周期图在 $\omega = 2\pi/20$ 和 $\omega = 2\pi/7$ 处出现峰值，主导周期 $T = 20$（Fisher 检验显著）。有限谐波回归：正则化最小二乘估计得到 $A_1 \approx 5$、$\omega_1 \approx 2\pi/20$、$A_2 \approx 2$、$\omega_2 \approx 2\pi/7$。贝叶斯更新后验集中在 $(\omega_1, \omega_2)$ 处。窗验证：$\mathcal{V}(K_{\text{Gauss}, \Delta=3}) \approx 0.95$，$\mathcal{V}(K_{\text{Cauchy}, \Delta=20}) \approx 0.88$，残差诊断 $\mathcal{D} \approx 0.1$。
\end{example}

\begin{theorem}[模型选择准则]\label{thm:ts_model_selection}
设 $\hat\sigma^2_k$ 为候选模型的极大似然方差估计，$p_k$ 为参数个数。则 AIC 与 BIC 准则为
\[
\operatorname{AIC}(k)=n\log\hat\sigma_k^2+2p_k,\qquad
\operatorname{BIC}(k)=n\log\hat\sigma_k^2+p_k\log n.
\]
在候选模型集合中，AIC 与 BIC 分别渐近选择使预测误差与后验概率最优的模型。
\end{theorem}

\begin{proof}
AIC 是 KL 散度与模型复杂度之和的渐近无偏估计；BIC 是后验模型概率的 Laplace 逼近。
\end{proof}

\begin{remark}[适用范围]\label{rem:ts_scope}
本章方法适用于具有平稳或局部平稳结构的离散时间序列，包括趋势提取、周期识别、拐点检测和短期预测。对于非平稳、长记忆或强非线性序列，应结合差分、局部建模或非线性模型使用。
\end{remark}

\begin{remark}[与 Kalman 滤波的接口]\label{rem:ts_kalman_interface}
对状态空间形式
\[
x_{t+1}=Fx_t+w_t,\qquad y_t=Hx_t+v_t,
\]
Kalman 滤波给出最优线性状态估计
\[
\hat x_{t|t}=\hat x_{t|t-1}+K_t(y_t-H\hat x_{t|t-1}),
\]
其中 $K_t$ 为 Kalman 增益。ARMA 模型可通过状态空间实现纳入 Kalman 框架，从而统一处理缺失值与在线更新。
\end{remark}

\section*{本章七主题对照表}

\begin{center}
\renewcommand{\arraystretch}{1.3}
\begin{tabular}{>{\centering\arraybackslash}p{2.2cm}>{\centering\arraybackslash}p{10.5cm}}
\toprule
主题 & 本章对应 \\
\midrule
平移 & 滞后算子 $B$（定义 \ref{def:ts_lag_operator}）生成离散平移群；窗函数卷积平滑 \\
缩放 & 窗宽 $\Delta$ 的尺度选择；保真度 $\mathcal{V}(K_\Delta)$（定义 \ref{def:ts_window_validation}） \\
旋转 & 瞬时极限角 $\theta_t$（定义 \ref{def:ts_limit_angle}）；相似度核（命题 \ref{prop:ts_turning_similarity}） \\
守恒 & 周期图能量守恒；Fisher 检验量归一化（命题 \ref{prop:ts_period_significance}） \\
对称 & 余弦对称（引理 \ref{lem:ts_turning_equivalence}）；方向反转（定理 \ref{thm:ts_refraction_relation}） \\
闭合 & 算法流程闭环（定义 \ref{def:ts_algorithm}）；白噪声残差 $\mathcal{D}\to 0$（命题 \ref{prop:ts_residual_diagnosis}） \\
周期 & 主导周期（定义 \ref{def:ts_period_detection}）；Fisher 检验（命题 \ref{prop:ts_period_significance}） \\
\bottomrule
\end{tabular}
\end{center}

\section*{本章思考题}

\begin{enumerate}[leftmargin=3em, label=\textbf{题 \arabic*}]
    \item 证明滞后算子的代数性质（性质 \ref{prop:ts_lag_algebra}）：
    \begin{enumerate}[label=(\roman*)]
        \item 线性性：$B(\alpha X_t+\beta Y_t)=\alpha BX_t+\beta BY_t$；
        \item 可交换性：$B^aB^b=B^{a+b}=B^bB^a$；
        \item 可逆性：当特征多项式 $\phi(z)$ 的根都在单位圆外时，$\phi(B)$ 可逆。
    \end{enumerate}
    请用序列空间上的线性算子语言完成证明，
    并说明为什么滞后算子是离散平移群 $\mathbb{Z}$ 的生成元。

    \item 证明差分算子的零空间刻画（引理 \ref{lem:ts_diff_kernel}）：
    $\ker\nabla^d$ 由 $\{1,t,t^2,\dots,t^{d-1}\}$ 张成。
    请用归纳法完成证明，并说明差分算子如何消除多项式趋势。

    \item 证明 Wold 分解定理（定理 \ref{thm:ts_wold}）：
    零均值弱平稳序列可唯一分解为
    \[
    X_t=\sum_{j=0}^\infty\psi_j\varepsilon_{t-j}+V_t,\quad
    \psi_0=1,\quad\sum_j\psi_j^2<\infty.
    \]
    请用投影定理完成证明，
    并说明 $V_t$ 为确定性成分的几何含义。

    \item 证明 Yule-Walker 方程（定理 \ref{thm:ts_yule_walker}）：
    对 AR($p$) 模型，
    \[
    \rho(h)=\phi_1\rho(h-1)+\cdots+\phi_p\rho(h-p),\quad h=1,\dots,p.
    \]
    请用 $\mathbb{E}[\varepsilon_tX_{t-h}]=0$（$h\ge1$）完成证明，
    并说明为什么 Yule-Walker 方程给出 AR 参数估计的矩方法。

    \item 证明平稳性的几何直观（性质 \ref{prop:ts_stationarity_geometry}）：
    AR(1) 模型 $X_t=\phi X_{t-1}+\varepsilon_t$ 平稳当且仅当 $|\phi|<1$，
    此时 $\gamma(h)=\sigma^2\phi^{|h|}/(1-\phi^2)$。
    请用特征根在单位圆外条件和自协方差递推完成证明，
    并说明 $|\phi|$ 越接近 $1$ 记忆越长的原因。

    \item 根据定义 \ref{def:ts_ar} 至定义 \ref{def:ts_arch_garch}，
    写出 AR($p$)、MA($q$)、ARMA($p,q$)、ARIMA($p,d,q$)、
    SARIMA、ARCH($q$)、GARCH($p,q$) 七种模型的定义式，
    并验证表 \ref{tab:ts_classical_correspondence} 中各行对应的
    ACF/PACF 特征、频域特征与算子对应。

    \item 证明拐点阈值与余弦相似度的等价性（引理 \ref{lem:ts_turning_equivalence}）：
    设 $|\Delta\theta|\in[0,\pi)$，$\tau\in(0,\pi)$，则
    \[
    |\Delta\theta|>\tau\Longleftrightarrow\cos(\Delta\theta)<\cos\tau.
    \]
    请用余弦函数在 $[0,\pi]$ 上的严格单调递减性与偶函数性质完成证明。

    \item 证明命题 \ref{prop:ts_turning_similarity} 中拐点处的三条相似度特征：
    \begin{enumerate}[label=(\roman*)]
        \item 余弦相似度 $\operatorname{Similarity}^{(1)}=\cos(\Delta\theta)<\cos\tau$；
        \item 正弦相似度 $\operatorname{Similarity}^{(2)}=\sin(\Delta\theta)$ 的符号给出转向；
        \item 倍角正弦相似度 $\operatorname{Similarity}^{(4)}=\sin(2\Delta\theta)$
              的绝对值在 $|\Delta\theta|<\pi/4$ 时随 $|\Delta\theta|$ 严格递增。
    \end{enumerate}
    请分别写出证明，并说明为什么正弦相似度的符号变化精确标记了拐点。

    \item 验证例 \ref{ex:ts_turning_example}：
    序列在 $t=50$ 处从上升转为下降，
    $\tilde{X}_{49}=10.2$，$\tilde{X}_{50}=10.5$，$\tilde{X}_{51}=10.1$。
    请计算：
    \begin{enumerate}[label=(\roman*)]
        \item 瞬时极限角 $\theta_{50}$ 与 $\theta_{51}$；
        \item 差角 $\Delta\theta$；
        \item 余弦相似度 $\operatorname{Similarity}^{(1)}$ 与正弦相似度 $\operatorname{Similarity}^{(2)}$；
        \item 判断在阈值 $\tau=30^\circ$ 下 $t=51$ 是否为拐点。
    \end{enumerate}

    \item 证明定理 \ref{thm:ts_refraction_relation}（转折率的半角分解与界）：
    设 $\theta_1,\theta_2\in(-\pi/2,\pi/2)$，转折率
    \[
    R_{\text{turn}}(t^*)=|\sin\theta_1-\sin\theta_2|,
    \]
    则
    \[
    R_{\text{turn}}=2\Bigl|\cos\frac{\theta_1+\theta_2}{2}\Bigr|
    \cdot\Bigl|\sin\frac{\theta_1-\theta_2}{2}\Bigr|,
    \]
    且 $0\le R_{\text{turn}}\le2$。
    请用和差化积公式完成证明，
    并验证三个特例：$R=0\iff\theta_1=\theta_2$；
    $R=2\iff\theta_2=-\theta_1=\pm\pi/2$；
    方向反转时 $R=2|\sin\theta_1|$。

    \item 证明命题 \ref{prop:ts_period_significance}（Fisher 检验）：
    设周期图在 $m=\lfloor(n-1)/2\rfloor$ 个正频率点上的取值为 $I_1,\dots,I_m$，
    Fisher 检验统计量
    \[
    g=\frac{\max_kI_k}{\sum_{k=1}^mI_k},
    \]
    显著性阈值为 $g_\alpha=1-(\alpha/m)^{1/(m-1)}$。
    请说明白噪声下周期图纵坐标独立同分布于指数分布，
    归一化比值服从 $\operatorname{Beta}(1,m-1)$ 分布，
    并用并集界完成证明。

    \item 证明命题 \ref{prop:ts_periodogram_distribution}（周期图渐近分布）：
    高斯白噪声下，周期图 $I_n(\omega_k)$ 渐近服从均值 $\sigma^2$ 的指数分布，
    且不同频率处渐近独立。
    请说明为什么 DFT 系数为复高斯随机变量。

    \item 证明命题 \ref{prop:ts_harmonic_estimate}（有限谐波模型参数的正则化最小二乘估计）：
    目标函数
    \[
    \mathcal{L}(\boldsymbol{\Theta})=\sum_{t=1}^n
    \Bigl(X_t-\operatorname{Re}\Bigl[\sum_kA_ke^{i(\omega_kt+\phi_k)}\Bigr]\Bigr)^2
    +\lambda\sum_k|A_k|,
    \]
    在适当的正则化条件下，估计量依概率收敛于真实参数。
    请说明 Wald 一致性条件的三个要点：
    参数空间紧致、真实参数唯一最小、损失函数逐点收敛。

    \item 证明定理 \ref{thm:ts_harmonic_asymptotic}（有限谐波估计的渐近正态性）：
    \[
    \sqrt{n}\bigl(\hat{\boldsymbol{\Theta}}_n-\boldsymbol{\Theta}_0\bigr)
    \xrightarrow{d}\mathcal{N}(0,\mathcal{I}(\boldsymbol{\Theta}_0)^{-1}).
    \]
    请用 Taylor 展开与 Slutsky 定理完成证明。

    \item 证明命题 \ref{prop:ts_bayes_convergence}（后验收敛性与周期稳定性）：
    在标准正则条件下，
    \[
    \sqrt{n}(\boldsymbol{\Theta}-\hat{\boldsymbol{\Theta}}_n)\mid X_{1:n}
    \xrightarrow{d}\mathcal{N}(0,\mathcal{I}(\boldsymbol{\Theta}_0)^{-1}).
    \]
    请说明后验分布在真实参数处集中、宽度为 $O(n^{-1/2})$，
    并指出如何用后验方差判断周期分量是否稳定。

    \item 证明命题 \ref{prop:ts_forgetting}（遗忘因子的自适应）：
    对非平稳时序，遗忘因子 $\gamma\in(0,1]$ 的定义为
    \[
    p_\gamma(\boldsymbol{\Theta}\mid X_{1:t})\propto
    p(X_t\mid\boldsymbol{\Theta})\cdot
    [p_\gamma(\boldsymbol{\Theta}\mid X_{1:t-1})]^\gamma,
    \]
    等价地
    \[
    \log p_\gamma(\boldsymbol{\Theta}\mid X_{1:t})\propto
    \sum_{s=1}^t\gamma^{t-s}\log p(X_s\mid\boldsymbol{\Theta}).
    \]
    请写出证明，并说明为什么 $\gamma\to1$ 时退化为标准贝叶斯更新，
    而 $\gamma<1$ 时旧观测权重指数衰减。

    \item 证明命题 \ref{prop:ts_window_tradeoff}（窗函数验证的拐点--周期权衡）：
    小窗口保真度高但噪声抑制不足；
    大窗口噪声抑制充分但高频信号被平滑；
    最优窗口在两者间平衡。
    请分别讨论 $\Delta\to0$ 与 $\Delta\to\infty$ 的极限行为，
    并说明最优窗口由最小化均方误差或交叉验证确定。

    \item 证明命题 \ref{prop:ts_residual_diagnosis}（残差的相似度诊断）：
    残差相似度诊断函数
    \[
    \mathcal{D}(\hat{\varepsilon})=\max_{h\ge1,\,d=1,\dots,6}
    \Bigl|\frac{1}{n-h}\sum_{t=h+1}^n
    S^{(d)}(\theta_t^{\hat{\varepsilon}},
    \theta_{t-h}^{\hat{\varepsilon}})\Bigr|.
    \]
    请说明理想白噪声残差的相位序列独立，
    差角族核的期望为零，由大数定律 $\mathcal{D}\to0$；
    并解释为什么 $\mathcal{D}$ 显著非零表明残差存在未捕获的周期结构。

    \item 证明定理 \ref{thm:ts_ljung_box}（Ljung-Box 残差白噪声检验）：
    残差自相关的检验统计量
    \[
    Q=n(n+2)\sum_{h=1}^m\frac{\hat\rho(h)^2}{n-h}
    \xrightarrow{d}\chi^2_{m-p-q}.
    \]
    请说明为什么在模型正确设定下 $Q$ 渐近服从卡方分布。

    \item 证明定理 \ref{thm:ts_model_selection}（模型选择准则）：
    AIC 与 BIC 准则
    \[
    \operatorname{AIC}(k)=n\log\hat\sigma_k^2+2p_k,
    \quad
    \operatorname{BIC}(k)=n\log\hat\sigma_k^2+p_k\log n.
    \]
    请说明 AIC 与 BIC 分别对应 KL 散度与后验模型概率的渐近近似。

    \item 验证例 \ref{ex:ts_complete_example}：
    模拟数据
    \[
    X_t=5\sin(2\pi t/20)+2\sin(2\pi t/7)+\varepsilon_t,
    \]
    其中 $\varepsilon_t\sim\mathcal{N}(0,1)$。
    请说明：
    \begin{enumerate}[label=(\roman*)]
        \item 拐点检测如何识别方向转折；
        \item 周期图在哪些频率处出现峰值；
        \item 有限谐波回归如何估计振幅与频率；
        \item 窗函数验证如何选择最优窗口。
    \end{enumerate}

    \item 证明定义 \ref{def:ts_algorithm} 中完整算法流程的闭环性：
    窗函数选择与平滑、瞬时极限角计算、拐点检测、转折率计算、
    周期确定、有限谐波回归、贝叶斯更新、窗函数验证
    八个步骤构成完整闭环。
    请说明为什么窗函数验证的结果可以反馈到第一步的窗函数选择。

    \item 综合题：结合滞后算子（定义 \ref{def:ts_lag_operator}）、
    瞬时极限角（定义 \ref{def:ts_limit_angle}）、
    相似度核（命题 \ref{prop:ts_turning_similarity}）、
    周期图（定义 \ref{def:ts_period_detection}）
    与残差诊断（命题 \ref{prop:ts_residual_diagnosis}），
    说明本章如何把“标量与矢量”章的相似度与趋势框架
    从三角函数层面扩充到离散时间序列层面。
    请分别指出：
    \begin{enumerate}[label=(\roman*)]
        \item 瞬时极限角如何对应向量层面的辐角；
        \item 拐点检测如何对应方向转折判定；
        \item 转折率如何对应趋势之差的度量；
        \item 残差诊断如何用相似度核检测未捕获周期。
    \end{enumerate}
\end{enumerate}

